\documentclass[11pt]{article}
\usepackage[margin=1in]{geometry}
\usepackage[T1]{fontenc}
\usepackage{lmodern,amsmath,amssymb,amsthm,mathtools,hyperref}
\setlength{\emergencystretch}{3em}
\newtheorem{theorem}{Theorem}[section]
\newtheorem{lemma}[theorem]{Lemma}
\newtheorem{corollary}[theorem]{Corollary}
\newtheorem{proposition}[theorem]{Proposition}
\newcommand{\tr}{\operatorname{tr}}
\title{Tensor stability, activation thresholds\\and limitations of dimension reduction}
\author{Joshua Mark\thanks{ORCID: \href{https://orcid.org/0009-0003-5905-2382}{0009-0003-5905-2382}. AI tools assisted with derivation, independent review, computation and exposition.}}
\date{20 November 2026 --- research draft; priority not established}
\begin{document}
\maketitle
\begin{abstract}
We give an elementary tangent-trace proof that the orthogonal matrices
whose every self-tensor power locally minimizes entrywise $p$-energy,
$p>2$, are exactly scaled real Hadamard matrices. The same argument
characterizes tensor-stable local maxima for $1<p<2$. On the complex
unitary group the characterization holds for every $p>0$, $p\ne2$,
including $p=1$. The real endpoint $p=1$ fails. The proof supplies
explicit finite-power necessary inequalities, a two-column curvature
certificate, and a dimension-independent quartic flatness estimate.
These are characterization and instability results, not a proof of
real Hadamard existence. Priority has not been established.
\end{abstract}

\section{Tensor-stable extrema over the real and complex fields}
Let $F_p(U)=\sum_{ij}|u_{ij}|^p$. A row-orthonormal matrix with $n$
columns is \emph{flat} if every entry has magnitude $n^{-1/2}$.
For a square real matrix this says precisely that $\sqrt nU$ is a
Hadamard matrix. Over $\mathbb C$, it says that $\sqrt nU$ is complex
Hadamard. We use the real Frobenius metric on the tangent spaces.

\begin{theorem}[Tensor stability]\label{thm:unified}
Let $n\ge2$ and let $U$ be square orthogonal or unitary.
\begin{enumerate}
\item Over $\mathbb R$, for every finite $p>2$, all positive
self-tensor powers of $U$ are local minima of $F_p$ if and only if
$U$ is flat. For $1<p<2$, replace minima by maxima.
\item Over $\mathbb C$, for every finite $p>2$, the same local-minimum
characterization holds. For every $0<p<2$, the local-maximum
characterization holds, including $p=1$.
\end{enumerate}
Every nonflat critical point in the indicated smooth domain has a
strict second-variation violation at some finite tensor power. For
$p<2$ this assertion assumes that its entries are nonzero; a matrix
with a zero entry is already excluded as a local maximum.
\end{theorem}

The orthogonal $p>2$ conclusion was established in the preceding
research pass by a factored direction. Here the proof needs neither
spectral assumptions on a critical matrix nor a factorization of the
Hessian witness. The real and complex row-Stiefel analogues also hold
by the same right-action computation below.

\begin{lemma}[Zeros and the strict moment comparison]\label{lem:moments}
For $0<p<2$, a row-orthonormal real or complex matrix with a zero entry
cannot locally maximize $F_p$ under right unitary or orthogonal
rotations. For a matrix without zero entries put
\[
 q=F_p(U),\quad r=\sum_{ij}|u_{ij}|^{p-2},\quad \rho=\frac{r}{nq}.
\]
For $p>2$ the same definitions allow zeros. Then $\rho<1$ when
$p>2$ and $U$ is nonflat, whereas $\rho>1$ when $0<p<2$ and $U$ is
nonflat. Equality occurs precisely at flat matrices.
\end{lemma}
\begin{proof}
If $u_{ij}=0$, choose $\ell$ with $u_{i\ell}\ne0$ and rotate columns
$j,\ell$ through angle $t$. In
$F_p(U(t))+F_p(U(-t))-2F_p(U)$, entries initially nonzero contribute
$O(t^2)$, since their linear terms cancel. Initially zero entries
contribute $2c|t|^p+o(|t|^p)$ with total $c>0$. This is positive for
small nonzero $t$ because $p<2$, contradicting a local maximum.
This argument also covers the nondifferentiable range $p\le1$.

For $d$ rows list the $dn$ squared magnitudes as $x_i$, so
$\sum x_i=d$, and set $a=(p-2)/2$. The identity
\[
 dnq-dr=\sum_{i<j}(x_i-x_j)(x_i^a-x_j^a)
\]
is strictly positive for $a>0$ and strictly negative for $a<0$
unless all $x_i$ agree. This proves the assertions.
\end{proof}

\section{The two tangent-trace identities}
The real identity is proved explicitly in
Theorem~\ref{thm:trace-test} below. For a square matrix, the trace over
a Frobenius-orthonormal skew basis is
\begin{equation}\label{eq:real-trace}
 \operatorname{Tr}_{\mathbb R}\operatorname{Hess}F_p(U)
 =\frac p2\big((p-1)r-(n+p-2)q\big).
\end{equation}
The formula here denotes the sum of second derivatives along
$Ue^{tX}$; it is valid without criticality. At a critical point it
is the tangent-Hessian trace. The normalized basis is
$(E_{ab}-E_{ba})/\sqrt2$ for $a<b$.

\begin{lemma}[Complex tangent trace]\label{lem:complex-trace}
For a row-orthonormal complex matrix with $n$ columns,
\begin{equation}\label{eq:complex-trace}
 \sum_X\frac{d^2}{dt^2}F_p(Ue^{tX})\bigg|_{t=0}
 =\frac p2\big(pr-(2n+p-2)q\big),
\end{equation}
where the sum is over the Frobenius-orthonormal anti-Hermitian basis
\[
 iE_{aa},\qquad (E_{ab}-E_{ba})/\sqrt2,\qquad
 i(E_{ab}+E_{ba})/\sqrt2\quad(a<b).
\]
For $p<2$ assume nonzero entries; for $p>2$ zeros are allowed.
\end{lemma}
\begin{proof}
For a fixed generator, write $z=u_{ij}$, $v=(UX)_{ij}$ and
$w=(UX^2)_{ij}$. At $z\ne0$, differentiation gives
\[
 \frac1p\frac{d^2}{dt^2}|(Ue^{tX})_{ij}|^p\bigg|_{t=0}
 =|z|^{p-2}|v|^2+(p-2)|z|^{p-4}
       \big(\Re(\bar zv)\big)^2+|z|^{p-2}\Re(\bar zw).
\]
At a zero entry for $p>2$ this second derivative is zero.
Summing over the stated basis gives, pointwise,
\[
 \sum_X |(UX)_{ij}|^2=1,\qquad
 \sum_X\big(\Re(\bar u_{ij}(UX)_{ij})\big)^2
       =\frac{|u_{ij}|^2(1-|u_{ij}|^2)}2,
 \qquad \sum_X X^2=-nI.
\]
For the middle identity, the diagonal imaginary generator contributes
zero; each real and imaginary off-diagonal pair contributes
$|u_{ij}|^2|u_{ia}|^2/2$. The first identity uses the same pairs and
the unit norm of row $i$. Substitution and summation over entries
give \eqref{eq:complex-trace}.
\end{proof}

\begin{proof}[Proof of Theorem~\ref{thm:unified}]
At $U^{\otimes k}$ the column number is $n^k$ and the two moments are
$q^k,r^k$. Dividing the real trace by $pn^kq^k/2$ gives
\[
 (p-1)\rho^k-1-\frac{p-2}{n^k}.
\]
For a nonflat matrix this is eventually negative for $p>2$, and
eventually positive for $1<p<2$. Dividing the complex trace by the
same positive quantity gives
\[
 p\rho^k-2-\frac{p-2}{n^k},
\]
which is eventually negative for $p>2$ and eventually positive for
every $0<p<2$. A negative trace has a negative summand, and a positive
trace has a positive summand. This contradicts the appropriate local
extremum. Lemma~\ref{lem:moments} handles zeros in the maximum case.
At critical points the exhibited second derivatives are Hessian
violations. Jensen inequality gives the converse: flat matrices are
global minima for $p>2$, global maxima for $0<p<2$, and remain flat
under tensor powers. The row-orthonormal argument is identical; the
right-rotation curves remain in its manifold.
\end{proof}

For example, sufficient strict-violation powers in the complex case
are $k>\log(p/2)/\log(1/\rho)$ for $p>2$ and
$k>\log(2/p)/\log\rho$ for $0<p<2$. These bounds need not be sharp.
The real maximum bound is $k>\log(1/(p-1))/\log\rho$.
Criticality of the seed propagates to tensor powers in the nonzero
domain: if $W=|U|^{p-2}\circ U$, criticality means $U^*W$ is
self-adjoint, and the corresponding tensor matrix is $(U^*W)^{\otimes k}$.

\section{Endpoints, prior work, and the existence gap}
At $p=2$ the energy is constant. The real endpoint $p=1$ fails:
for every $n>4$, $U=2J_n/n-I_n$ is nonflat, orthogonal, and all its
self-tensor powers are strict local maxima of $F_1$ on the real
orthogonal group. Indeed its sign matrix is $A=J_n-2I_n$, and
\[
 U^TA=(1-4/n)J_n+2I_n\succ0.
\]
Near $U$ the signs are fixed, so the Hessian is
$-\operatorname{tr}(U^TA X^TX)<0$ for nonzero real skew $X$.
Tensor powers preserve positivity of this critical matrix.
This endpoint behavior is consistent with the established real
almost-Hadamard theory~\cite{BN,BNZ}.

The complex almost-Hadamard conjecture of Banica and Nechita asks
whether a \emph{single} local maximum of $F_1$ on $U(n)$ must be
flat~\cite{Complex}. Our complex theorem assumes local maximality at
\emph{every} self-tensor power. It therefore does not settle that
stronger one-copy conjecture. In particular, tensor closure of a
class of balanced critical points does not establish tensor closure
of local maxima. Targeted searches did not establish bibliographic
priority for the tensor characterization; no novelty is claimed.

For real Hadamard existence the missing implication is different:
compactness supplies a local or global minimum in each dimension,
but it does not supply one whose every tensor power is locally
minimizing. Lifting and improving a nonflat point changes the
dimension. Neither this theorem nor its quantitative corollaries
extract an original-dimensional Hadamard matrix.
\section{A trace certificate and quantitative finite-power stability}
This section gives a second proof of tensor destabilization, using the
trace of the entire tangent Hessian. It produces a two-column rotation
in the lifted dimension and a bound on the energy of any matrix whose
specified tensor power is locally minimizing. Neither conclusion
supplies a descent in the original dimension.

Let $U\in\mathbb R^{d\times n}$ have orthonormal rows, $1\le d\le n$,
$n\ge2$, and let $p>2$. Write
\[
 q=\sum_{ij}|u_{ij}|^p,\qquad r=\sum_{ij}|u_{ij}|^{p-2},
 \qquad \rho=\frac{r}{nq}.
\]
By strict Chebyshev inequality, $0<\rho\le1$, with equality exactly
when every entry has magnitude $n^{-1/2}$.

\begin{theorem}[Finite-power trace test]\label{thm:trace-test}
If $U^{\otimes k}$ is a local minimum of the entrywise $p$-energy on
the row-Stiefel manifold, then
\begin{equation}\label{eq:trace-test}
 (p-1)\rho^k\ \ge\ 1+\frac{p-2}{n^k}.
\end{equation}
In particular, a nonflat $U$ has a strictly negative second variation
along a single two-column rotation of $U^{\otimes k}$ whenever
$k>\log(p-1)/\log(1/\rho)$. The second-variation conclusion does not
require that $U$ be critical.
\end{theorem}

\begin{proof}
For any real row-orthonormal $V$ with $N$ columns, let
$Q=\sum |v_{ij}|^p$ and $R=\sum |v_{ij}|^{p-2}$. For
$X_{ab}=E_{ab}-E_{ba}$, direct differentiation of
$F_p(Ve^{tX_{ab}})$ gives
\begin{align*}
 \frac1p\frac{d^2}{dt^2}F_p(Ve^{tX_{ab}})\bigg|_{t=0}
 ={}&(p-1)\sum_i\bigl(|v_{ia}|^{p-2}v_{ib}^2
               +|v_{ib}|^{p-2}v_{ia}^2\bigr)\\
 &-\sum_i\bigl(|v_{ia}|^p+|v_{ib}|^p\bigr).
\end{align*}
Summing over $a<b$ and using $\sum_bv_{ib}^2=1$ yields
\begin{equation}\label{eq:curvature-trace}
 \frac1p\sum_{a<b}\frac{d^2}{dt^2}F_p(Ve^{tX_{ab}})\bigg|_{t=0}
   =(p-1)R-(N+p-2)Q.
\end{equation}
For $V=U^{\otimes k}$ we have $N=n^k$, $Q=q^k$, $R=r^k$.
A local minimum has nonnegative second derivative on every such
curve, which gives \eqref{eq:trace-test}. If
$(p-1)\rho^k<1$, the displayed sum is strictly negative; at least one
of its summands is strictly negative. Smoothness at zero is valid
because $p>2$.
\end{proof}

The column pair can be checked without forming the tensor matrix.
Define $c_a=\sum_i|u_{ia}|^p$ and
$D_{ab}=\sum_i|u_{ia}|^{p-2}u_{ib}^2$. For distinct column words
$\mathbf a,\mathbf b\in\{1,\ldots,n\}^k$, the second derivative divided
by $p$ is exactly
\[
 (p-1)\left(\prod_{j=1}^kD_{a_jb_j}
             +\prod_{j=1}^kD_{b_ja_j}\right)
       -\prod_{j=1}^kc_{a_j}-\prod_{j=1}^kc_{b_j}.
\]
This is an explicit certificate formula, not a claim of a polynomial
time search algorithm. A two-column rotation of a tensor matrix
usually does not factor as rotations of its original factors.

\begin{corollary}[Dimension-independent quartic stability bound]
\label{cor:quartic-rate}
Let $e=nF_4(U)/d\ge1$. If $U^{\otimes k}$ is a local minimum, then
\[
 e\le \left(\frac{3}{1+2/n^k}\right)^{1/k}<3^{1/k}.
\]
Furthermore
\[
 \sum_{ij}\left(u_{ij}^2-\frac1n\right)^2
      =\frac dn(e-1),\qquad
 \frac1d\min_{H\in\{\pm1\}^{d\times n}}
       \left\|U-\frac H{\sqrt n}\right\|_F^2
      \le e-1<3^{1/k}-1.
\]
Thus the normalized squared distance to a flat sign array is
$O(1/k)$, uniformly in the matrix dimensions.
\end{corollary}

\begin{proof}
At $p=4$ we have $r=d$, hence $\rho=1/e$; apply
\eqref{eq:trace-test}. The first distance identity follows by expanding
the square and using $\sum u_{ij}^2=d$. Choose each sign of $H$ to
match $u_{ij}$ (arbitrarily at zero) and use
\[
 (|u_{ij}|-n^{-1/2})^2
 =\frac{(u_{ij}^2-n^{-1})^2}{(|u_{ij}|+n^{-1/2})^2}
 \le n(u_{ij}^2-n^{-1})^2.
\]
Finally $3^{1/k}-1\sim(\log3)/k$.
\end{proof}

The closest flat sign array in this estimate need not have orthogonal
rows. The estimate therefore proves neither a nearby Hadamard matrix
nor Hadamard existence. It quantifies a necessary condition on local
minima in tensor powers.

\subsection{Complex flatness and collision entropy}
For complex row-orthonormal $U$, put again $e=nF_4(U)/d$. The complex
trace identity gives the slightly stronger necessary condition
\[
 e\le\left(\frac{2}{1+1/n^k}\right)^{1/k}<2^{1/k}
\]
when $U^{\otimes k}$ is locally minimizing. The same variance identity
holds with $u_{ij}^2$ replaced by $|u_{ij}|^2$; replacing signs by
arbitrary unit complex phases gives the corresponding distance bound.
This follows by choosing the phase of each entry to match $u_{ij}$.

There is a direct information-theoretic interpretation of both bounds.
Define a probability distribution on the $dn$ entry positions by
$P_{ij}=|u_{ij}|^2/d$, and let $Q_{ij}=1/(dn)$ be uniform. Its
order-two R\'enyi divergence from uniform, here defined explicitly, is
\[
 D_2(P\Vert Q)=\log\sum_{ij}\frac{P_{ij}^2}{Q_{ij}}
             =\log e.
\]
The entry distribution of $U^{\otimes k}$ is $P^{\otimes k}$, so
$D_2(P^{\otimes k}\Vert Q^{\otimes k})=kD_2(P\Vert Q)$.
Consequently finite-power local minimality forces
\[
 D_2(P\Vert Q)<\frac{\log3}{k}\quad\text{over }\mathbb R,
 \qquad
 D_2(P\Vert Q)<\frac{\log2}{k}\quad\text{over }\mathbb C.
\]
This expresses the curvature test as an entropy-deficit bound. It
does not assert that an arbitrary nearly uniform entry distribution
can be realized by a flat orthogonal matrix of the original size.

\section{Exact activation thresholds and delayed instability}

Put $F_p(U)=\sum_{ij}|u_{ij}|^p$ on $O(n)$, where $p>2$.
A negative tangent Hessian means negative second derivative along
$Ue^{tX}$ for a real skew-symmetric $X$. This note sharpens the finite
tensor activation theorem with exact examples; it does not establish
Hadamard existence or claim priority.

\begin{theorem}\label{activation:main}
Let $A=2J_5/5-I_5$, $s=p-2$, and
\[
 b=2/5,\qquad c=(3/5)^{p-1}+(2/5)^{p-1},\qquad
 R(p)=c/b^s=\tfrac35(\tfrac32)^s+\tfrac25.
\]
For every positive integer $k$, the least eigenvalue of the tangent
Hessian of $F_p$ at $A^{\otimes k}$, with the Frobenius metric on
skew generators, is exactly
\[
 \lambda_k=p\big((p-1)b^{ks}-c^k\big).
\]
Consequently it is a strict local minimum if $R(p)^k<p-1$, has positive
semidefinite singular Hessian if equality holds, and has a negative
tangent direction if $R(p)^k>p-1$. The first negative-Hessian power is
\[
 k_-(p)=\left\lfloor\frac{\log(p-1)}{\log R(p)}\right\rfloor+1.
\]
The equality case does not assert local minimality.
\end{theorem}
\begin{proof}
Write $d=(3/5)^{p-1}$ and $e=(2/5)^{p-1}$. The entrywise gradient
matrix is $W=eJ_5-(d+e)I_5$, so $M=A^TW$ is symmetric and $A$ is
critical. On $\mathbf1^\perp$, $M$ has eigenvalue $c=d+e$, while on
$\mathbf1$ it has eigenvalue $h=-d+4e$. For $p>2$,
$3e<2d$, hence $-c<h<c$. Thus $\|M\|_{\rm op}=c$.
Tensor multiplicativity gives criticality and $M_k=M^{\otimes k}$.
Every entry of $A^{\otimes k}$ has magnitude at least $b^k$.
For any skew $X$, orthogonality and the operator norm bound imply
\[
 T_k(X):=\sum_{ij}|(A^{\otimes k})_{ij}|^s
                 ((A^{\otimes k}X)_{ij})^2\ge b^{ks}\|X\|_F^2,
 \qquad
 S_k(X):=\tr(M_kX^TX)\le c^k\|X\|_F^2.
\]
Here the last bound remains valid if $M_k$ has negative eigenvalues.
The derivative identity $F_p''=p[(p-1)T_k-S_k]$ gives the claimed
lower bound for the least Hessian eigenvalue.

To attain it, take a skew matrix $Y$ supported on the oriented triangle
$0\to1\to2\to0$, with entries $+1$ on these three edges and $-1$
on reversed edges. Take symmetric $Z$ supported on the undirected
cycle $0,1,2,3,0$, with successive edge values $+1,-1,+1,-1$.
Both have zero diagonal and zero row and column sums. Consequently
$AY=-Y$, $AZ=-Z$, and $MY=cY$, $MZ=cZ$.
For $X=Y\otimes Z^{\otimes(k-1)}$, each nonzero entry has distinct
row and column indices in every coordinate. Thus it encounters only
magnitude-$b^k$ entries of $A^{\otimes k}$ in $T_k$, and
\[
 T_k(X)=b^{ks}\|X\|_F^2,\qquad S_k(X)=c^k\|X\|_F^2.
\]
The matrix $X$ is nonzero and skew, proving sharpness. Positive
definite Hessian at a critical point yields a strict local minimum.
\end{proof}

\begin{corollary}[Two copies do not suffice]
For $p=3$, $A$ and $A^{\otimes2}$ are strict local minima, while
$A^{\otimes3}$ has a negative tangent direction. For $p=201/100$,
the first four tensor powers are strict local minima, while the fifth
has a negative tangent direction.
\end{corollary}
\begin{proof}
At $p=3$, $R=13/10$ and $R^2<2<R^3$.
For $p=201/100$, exact integer exponentiation verifies
\[
 (1004062/10^6)^{100}<3/2<(1004063/10^6)^{100}.
\]
Put $a=\frac35(1004062/10^6)+\frac25$ and
$z=\frac35(1004063/10^6)+\frac25$. Exact rational arithmetic gives
$z^4<101/100<a^5$. Hence $R^4<p-1<R^5$.
\end{proof}

The example at $p=3$ refutes universal two-copy negative-Hessian
activation, including when the seed is a strict local minimum.
The first-power dimension here is five, so it is not itself an
admissible Hadamard order. Tensoring with a normalized Hadamard matrix
also preserves this exact Hessian classification: the minimum entry
weight and the operator norm of the critical matrix both acquire
the same positive factor. An attaining generator is the generator
above tensored with the identity on the flat factor. Thus taking
$A\otimes(H_4/2)$ gives the same examples in initial dimension twenty.
For its $k$th self-power the flat factor has order $4^k$; the Hessian
least eigenvalue is multiplied by $(4^k)^{-(p-2)/2}$.

All activation indices in this specific family are at most five.
Indeed convexity of $\log(\frac35e^{s\log(3/2)}+\frac25)$ gives
$\log R\ge\frac35s\log(3/2)$, whereas $\log(1+s)<s$.
Since $3\log(3/2)>1$, it follows that $5\log R>\log(p-1)$.
Moreover the threshold tends to $5/(3\log(3/2))\simeq4.1105$ as
$p\downarrow2$. This proves that five is attained, but does not
settle whether other critical points require arbitrarily many copies.

\section{Factored directions for comparison}
For the following refinement, let $U\in O(n)$ be critical, put
$W=|U|^{p-2}\circ U$ and $M=U^TW=M^T$, and define
\[
 T(V)=\sum_{ij}|u_{ij}|^{p-2}(UV)_{ij}^2,
 \qquad S(V)=\tr(MV^TV).
\]
The second variation in a skew generator is $p((p-1)T-S)$.
The symmetric and skew orthonormal basis trace sums are respectively
$((r+q)/2,(n+1)q/2)$ and $((r-q)/2,(n-1)q/2)$ for $(T,S)$.
Since $r<nq$, the ratio-selection argument selects symmetric $Z$ and
skew $Y$ with positive $S$ and ratios less than one. This remains
valid when some basis $S$-values are negative: if every positive-$S$
ratio exceeded the quotient of the sums, summing $T-cS$ would be
strictly positive, a contradiction (the nonpositive-$S$ terms are
nonnegative). Tensor products factor both quadratic forms. This is
the factored-direction proof preserved in the previous research pass.
\section{Interaction-only descent and an obstruction to dimension reduction}
\label{sec:interaction-descent}

Throughout this section $p>2$, $U\in O(n)$ is critical for
$F_p(U)=\sum_{ij}|u_{ij}|^p$, and $n\ge2$. Use the quantities
$q,r,T,S,M$ from the tensor-stability theorem. The following refinement
locates the destabilizing direction entirely outside the tangent space
of independently varying tensor factors.

\begin{proposition}[Trace-free activation]
If $U$ is nonflat, there are a skew-symmetric $Y$ and a symmetric
trace-zero $Z$ with $S(Y),S(Z)>0$ such that
\[
 \alpha=\frac{T(Y)}{S(Y)}\le\frac{r-q}{(n-1)q}<1,
 \qquad
 \beta=\frac{T(Z)}{S(Z)}\le
 b_0:=\frac{nr+(n-2)q}{(n+2)(n-1)q}<1.
\]
Moreover $0<b_0<b=(r+q)/((n+1)q)$. For any integer $k\ge2$ satisfying
$(p-1)\alpha\beta^{k-1}<1$, the generator
$X=Y\otimes Z^{\otimes(k-1)}$ gives negative tangent Hessian at
$U^{\otimes k}$. One sufficient choice is
\[
 k=2+\left\lceil\frac{\log(p-1)}{\log(1/b_0)}\right\rceil.
\]
\end{proposition}
\begin{proof}
In the symmetric matrices, the scalar direction $I_n/\sqrt n$ is
orthogonal to the trace-zero subspace. Its $T$- and $S$-values both
equal $q/n$. Subtracting these from the symmetric trace sums gives,
for any Frobenius-orthonormal basis $\mathcal B_0$ of the symmetric
trace-zero matrices,
\[
 \sum_{Z\in\mathcal B_0}T(Z)=\frac{r+q}{2}-\frac qn,
 \qquad
 \sum_{Z\in\mathcal B_0}S(Z)=\frac{(n+1)q}{2}-\frac qn.
\]
The second sum is positive. The ratio-selection argument therefore
produces the asserted $Z$; the skew selection is unchanged.
The strict inequality $r<nq$ gives $b_0<1$, and direct subtraction gives
\[
 b-b_0=\frac{2(nq-r)}{(n+1)(n+2)(n-1)q}>0.
\]
The numerator of $b_0$ is positive, including $n=2$, because $r>0$.
The tensor Hessian factors as
$pS(Y)S(Z)^{k-1}((p-1)\alpha\beta^{k-1}-1)$.
This proves both negativity and the sufficient bound.
\end{proof}

\begin{proposition}[The descent is invisible to one-factor tangents]
Let $k\ge2$, and let $X$ be a trace-free activating generator as above.
Every one-factor partial trace of $X$ vanishes. Equivalently, $X$ is
Frobenius-orthogonal to
\[
 \mathcal L_k=\left\{
 \sum_{a=1}^k I^{\otimes(a-1)}\otimes A_a\otimes
 I^{\otimes(k-a)}: A_a^{\mathsf T}=-A_a\right\}.
\]
This is the relative tangent space of the manifold of independent
tensor factors at $U^{\otimes k}$.
\end{proposition}
\begin{proof}
For elementary tensors, the partial trace retaining factor $a$ is
the $a$th factor multiplied by the traces of all the others.
For $a=1$, it contains $(\operatorname{tr}Z)^{k-1}=0$.
For $a>1$, it contains $\operatorname{tr}Y=0$. Partial trace is adjoint
to insertion of identities, proving the orthogonality statement.
Differentiating $\bigotimes_a(Ue^{tA_a})$ gives
$U^{\otimes k}\sum_a I^{\otimes(a-1)}\otimes A_a\otimes
I^{\otimes(k-a)}$. Left multiplication by the orthogonal
$U^{\otimes k}$ preserves the Frobenius inner product.
\end{proof}

\begin{proposition}[Independent factors cannot realize the lifted gain]
Suppose $U$ is a local minimum of $F_p$ on $O(n)$ and put $q=F_p(U)$.
Then $U^{\otimes k}$ is a local minimum of $F_p$ restricted to the
manifold of independent tensor factors, for every $k$. At a critical
$U$, its Hessian on relative product tangents is
\[
 \operatorname{Hess}F_p\big|_{\mathcal L_k}
 \left[\sum_a I^{\otimes(a-1)}\otimes A_a\otimes
 I^{\otimes(k-a)}\right]
 =q^{k-1}\sum_a \operatorname{Hess}F_p(U)[A_a].
\]
If $U$ is nonflat, choose the interaction-only negative direction $X$
above and let $W(t)=U^{\otimes k}e^{tX}$. For the local nearest-point
projection $P$ onto the product manifold,
\[
 F_p(W(t))=q^k-c t^2+o(t^2)\quad(c>0),\qquad
 q^k\le F_p(P(W(t)))=q^k+O(t^4).
\]
Thus nearest-product projection loses the entire quadratic improvement.
\end{proposition}
\begin{proof}
Entrywise energy is exactly multiplicative:
$F_p(\bigotimes_a V_a)=\prod_a F_p(V_a)$.
Nearby factors have energy at least $q$, proving restricted minimality.
Twice differentiating the product shows the stated Hessian: all cross
terms vanish by criticality. The local product map is an immersion,
since its factor tangent summands are pairwise orthogonal (each $A_a$
has trace zero); it is locally an embedded manifold. Its Euclidean
nearest-point projection is smooth near the base point, and its
differential is orthogonal projection onto the tangent space.
The preceding proposition gives
$P(W(t))=U^{\otimes k}+O(t^2)$. Restricted criticality and a bounded
second derivative give the $O(t^4)$ energy estimate, while restricted
minimality gives its sign. Taylor's theorem in the full orthogonal
group gives the negative quadratic coefficient for $W$.
\end{proof}

This is a precise limitation on one natural extraction procedure, not
a theorem excluding every nonlinear dimension-reduction argument.
The partial traces here are traces of the \emph{relative tangent
generator} $X$; no assertion about reduced physical states is needed.

\begin{proposition}[An exactly stationary partial-trace extraction]
 Let $k\ge2$ and $X=Y\otimes Z^{\otimes(k-1)}$, where $Y$ is real skew-symmetric,
$Z$ is real symmetric, and the eigenvalues of $Z$ occur in opposite
pairs, allowing zeros. Set
\[
 A_a(t)=n^{-(k-1)}\operatorname{Tr}_{\widehat a}(e^{tX}),
 \qquad a=1,\ldots,k,
\]
where the trace retains factor $a$. For all sufficiently small real
$t$, every $A_a(t)$ is symmetric positive definite. Consequently the
orthogonal polar factor of every $A_a(t)$ is exactly $I_n$.
Thus extracting $U\operatorname{polar}(A_a(t))$ from a lifted curve
returns the original $U$ in every factor, exactly rather than only
to first order.
\end{proposition}
\begin{proof}
Use the convergent power series for the exponential. In the partial
trace retaining factor $1$, the coefficient of $t^m/m!$ is
$Y^m(\operatorname{tr}Z^m)^{k-1}$; for odd $m$ this vanishes by the
paired spectrum of $Z$. All even powers of $Y$ are symmetric.
For a retained factor $a>1$, the coefficient is
$Z^m\operatorname{tr}(Y^m)(\operatorname{tr}Z^m)^{k-2}$.
Odd powers of a real skew-symmetric matrix have trace zero, and
all powers of $Z$ are symmetric. Thus all $A_a(t)$ are symmetric.
Since $A_a(0)=I_n$, continuity makes them positive definite near zero.
The orthogonal polar factor of a positive definite matrix is $I_n$.
\end{proof}

This obstruction occurs for a genuinely decreasing curve, not just an
arbitrary interaction. Take $p=4$, $U=2J_5/5-I_5$, and let $Y$ be the
oriented triangle on indices $0,1,2$, with nonzero entries
$Y_{01}=Y_{12}=Y_{20}=1$ and their skew counterparts. Let $Z$ be the
symmetric alternating four-cycle with
$Z_{01}=1,Z_{12}=-1,Z_{23}=1,Z_{30}=-1$ and their symmetric counterparts.
Then $Z$ has paired spectrum: conjugation by
$D=\operatorname{diag}(1,-1,1,-1,1)$ takes $Z$ to $-Z$.
The exact tangent Hessian at $U\otimes U$ in direction $Y\otimes Z$
is $-192/625$. Hence this full energy strictly decreases for sufficiently
small nonzero $t$, whereas the partial-trace polar extraction returns
$U$ exactly. To check the base minimum directly,
$M=\frac7{25}I_5-\frac6{125}J_5$ and every $u_{ij}^2\ge4/25$.
Thus $T(A)\ge(4/25)\|A\|_F^2$ and
$S(A)\le(7/25)\|A\|_F^2$, giving
$4(3T(A)-S(A))\ge(4/5)\|A\|_F^2$ for every skew $A$.
This proves that $U$ is a strict local minimum.

\subsection{An exact example where flatness does not descend}
Even an exactly flat tensor-square dimension need not permit a flat
matrix in the original dimension. There is no Hadamard matrix of order
$6$: normalizing one row to $+1$, orthogonality forces every other row
to have three positive entries; two such rows cannot have zero inner
product, since their positive-support intersection would have size
$6/4$, which is not an integer.

For completeness, an entirely explicit Hadamard matrix of order
$36=6^2$ can be built as follows. Let indices of $C$ be
$\{\infty\}\cup\mathbb F_{17}$, set its diagonal to zero,
$C_{\infty a}=C_{a\infty}=1$, and
$C_{ab}=\chi(a-b)$ for distinct finite indices, where $\chi$ is the
quadratic character modulo $17$. Then $C$ is symmetric and
$C^2=17I_{18}$. These are finite identities, also checked directly in
the accompanying exact verifier. Consequently
\[
 H=\begin{pmatrix}C+I&C-I\\ C-I&-C-I\end{pmatrix}
 \quad\text{satisfies}\quad HH^{\mathsf T}=36I_{36}.
\]
Indeed its diagonal Gram blocks are
$(C+I)^2+(C-I)^2=36I$, and its off-diagonal blocks vanish by
commutation. Thus a dimension-independent rule taking a flat
$n^2$-dimensional orthogonal matrix to a flat $n$-dimensional one is
impossible. The example uses a forbidden base order and does not
settle what can be extracted at orders divisible by four.

\begin{thebibliography}{9}
\bibitem{BN} T. Banica and I. Nechita, Almost Hadamard matrices:
the case of arbitrary exponents, Discrete Applied Mathematics 161
(2013), 2367--2379. \url{https://arxiv.org/abs/1211.2669}.
\bibitem{BNZ} T. Banica, I. Nechita and K. \.{Z}yczkowski,
Almost Hadamard matrices: general theory and examples,
Open Systems and Information Dynamics 19 (2012).
\url{https://arxiv.org/abs/1202.2025}.
\bibitem{Complex} T. Banica and I. Nechita,
Almost Hadamard matrices with complex entries (2017).
\url{https://arxiv.org/abs/1702.02220}.
\end{thebibliography}
\end{document}
